% --- Archon physics formalization source begin ---
% archon:physics
% archon:covers PhyXMiniProblems/problem_phyx_mini_0463.lean
% archon:source-report reports/phyx_mini/problem_phyx_mini_0463.source.json

\chapter{Physics problem phyx\_mini\_0463}
\label{ch:PhyXMiniProblems_problem_phyx_mini_0463}

\paragraph{Problem source.}
\#\# Physical scenario

Ten kilograms of water in a piston/cylinder arrangement exists as saturated liquid/vapor at \$100 \textbackslash{}, \textbackslash{}text\{kPa\}\$, with a quality of \$50\textbackslash{}\%\$. The system is now heated so that the volume triples. The mass of the piston is such that a cylinder pressure of \$200 \textbackslash{}, \textbackslash{}text\{kPa\}\$ will float it, as in figure.

\#\# Question

Find the heat transfer in the process.

\#\# Answer choices

- A: A: 20882.7 kJ

- B: B: 37187.9 kJ

- C: C: 3390 kJ

- D: D: 25961 kJ

\#\# Figure caption (auxiliary; use the image as primary evidence)

The image depicts a vertical cylindrical container partially filled with water, labeled "H₂O." There is a movable piston inside the container above the water. 

Text Description:
- The pressure above the piston is denoted as "P₀."
- There is a downward arrow labeled "g," indicating the direction of gravitational force.

Components and Relationships:
- The gray container holds the blue-colored water (H₂O).
- The piston separates the water from the space above it at pressure P₀.
- The piston is shown as a rectangular block that can move vertically within the container walls.
- An arrow pointing downward next to the container represents gravity (g) acting on the system. 

Overall, the diagram illustrates a basic setup of water under a piston in a cylindrical container influenced by gravity.

\#\# Dataset metadata

Laws of Thermodynamics; Physical Model Grounding Reasoning, Multi-Formula Reasoning

\paragraph{Recorded answer/context.}
D: D: 25961 kJ

\paragraph{Figure/image path.}
/root/proposal\_for\_physic/science-mango/phyx\_mini\_run/phyx\_data/test\_image/463.png

\paragraph{Formalization target.}
create a compiling Lean file with sorry bodies at `PhyXMiniProblems/problem\_phyx\_mini\_0463.lean`. The Lean declarations must preserve the physical quantities, dimensions or dimensional roles, figure labels, governing-law hypotheses, and final relation expressed by this problem.
Use Mathlib/Physlib names found through LeanExplore where available. If a physics API is missing, introduce faithful local abstractions rather than scalar placeholder aliases.

\begin{theorem}[Physics formalization target]
\label{thm:physics:phyx_mini_0463:target}
\lean{PhyXMiniProblems.ProblemPhyXMini0463.problem_phyx_mini_0463}
\uses{def:physics:phyx-mini-0463:phyxminiproblems-problemphyxmini0463-energyinkilojoules, def:physics:phyx-mini-0463:phyxminiproblems-problemphyxmini0463-stoppedpistonwatersetup, def:physics:phyx-mini-0463:phyxminiproblems-problemphyxmini0463-matchesstoppedpistonwaterscenario, def:physics:phyx-mini-0463:phyxminiproblems-problemphyxmini0463-matchesproblemreadouts, def:physics:phyx-mini-0463:phyxminiproblems-problemphyxmini0463-matchessuppliedstoppedpistonfigure, def:physics:phyx-mini-0463:phyxminiproblems-problemphyxmini0463-hasphysicalstoppedpistonparameters, def:physics:phyx-mini-0463:phyxminiproblems-problemphyxmini0463-satisfieswaterpropertyrelations, def:physics:phyx-mini-0463:phyxminiproblems-problemphyxmini0463-matchesexercisewaterpropertydata, def:physics:phyx-mini-0463:phyxminiproblems-problemphyxmini0463-satisfiesclosedwatermassvolumelaw, def:physics:phyx-mini-0463:phyxminiproblems-problemphyxmini0463-satisfiespistonloadbalance, def:physics:phyx-mini-0463:phyxminiproblems-problemphyxmini0463-satisfiesstoppedthenconstantpressurepath, def:physics:phyx-mini-0463:phyxminiproblems-problemphyxmini0463-satisfiestwostageboundaryworklaw, def:physics:phyx-mini-0463:phyxminiproblems-problemphyxmini0463-satisfiesclosedsystemfirstlaw, def:physics:phyx-mini-0463:phyxminiproblems-problemphyxmini0463-isuniqueclosestreportedheattransferchoice}
The assigned autoformalize agent should translate this physics problem into Lean declarations in the covered file, with theorem and lemma proof bodies written as `by sorry`.
\end{theorem}
\begin{proof}
This is an autoformalization task, not a proof task. Produce statements that can later be proved by the physics prover without weakening the physical model.
\end{proof}

% --- Archon declaration topology begin ---
\section{Declaration topology}
\begin{definition}[volume Dimension]
  \label{def:physics:phyx-mini-0463:phyxminiproblems-problemphyxmini0463-volumedimension}
  \lean{PhyXMiniProblems.ProblemPhyXMini0463.volumeDimension}
  The physical dimension of volume, L³
\end{definition}
\begin{proof}
  This is the declared model definition of volume Dimension.
\end{proof}
\begin{definition}[specific Volume Dimension]
  \label{def:physics:phyx-mini-0463:phyxminiproblems-problemphyxmini0463-specificvolumedimension}
  \lean{PhyXMiniProblems.ProblemPhyXMini0463.specificVolumeDimension}
  \uses{def:physics:phyx-mini-0463:phyxminiproblems-problemphyxmini0463-volumedimension}
  The physical dimension of specific volume, L³ M⁻¹
\end{definition}
\begin{proof}
  This is the declared model definition of specific Volume Dimension.
\end{proof}
\begin{definition}[specific Energy Dimension]
  \label{def:physics:phyx-mini-0463:phyxminiproblems-problemphyxmini0463-specificenergydimension}
  \lean{PhyXMiniProblems.ProblemPhyXMini0463.specificEnergyDimension}
  The physical dimension of specific energy, L² T⁻²
\end{definition}
\begin{proof}
  This is the declared model definition of specific Energy Dimension.
\end{proof}
\begin{definition}[acceleration Dimension]
  \label{def:physics:phyx-mini-0463:phyxminiproblems-problemphyxmini0463-accelerationdimension}
  \lean{PhyXMiniProblems.ProblemPhyXMini0463.accelerationDimension}
  The physical dimension of acceleration, L T⁻²
\end{definition}
\begin{proof}
  This is the declared model definition of acceleration Dimension.
\end{proof}
\begin{definition}[force Dimension]
  \label{def:physics:phyx-mini-0463:phyxminiproblems-problemphyxmini0463-forcedimension}
  \lean{PhyXMiniProblems.ProblemPhyXMini0463.forceDimension}
  \uses{def:physics:phyx-mini-0463:phyxminiproblems-problemphyxmini0463-accelerationdimension}
  The physical dimension of force, M L T⁻²
\end{definition}
\begin{proof}
  This is the declared model definition of force Dimension.
\end{proof}
\begin{definition}[Mass Quantity]
  \label{def:physics:phyx-mini-0463:phyxminiproblems-problemphyxmini0463-massquantity}
  \lean{PhyXMiniProblems.ProblemPhyXMini0463.MassQuantity}
  A nonnegative, unit-independent physical mass
\end{definition}
\begin{proof}
  This is the declared model definition of Mass Quantity.
\end{proof}
\begin{definition}[Volume Quantity]
  \label{def:physics:phyx-mini-0463:phyxminiproblems-problemphyxmini0463-volumequantity}
  \lean{PhyXMiniProblems.ProblemPhyXMini0463.VolumeQuantity}
  \uses{def:physics:phyx-mini-0463:phyxminiproblems-problemphyxmini0463-volumedimension}
  A nonnegative, unit-independent physical volume
\end{definition}
\begin{proof}
  This is the declared model definition of Volume Quantity.
\end{proof}
\begin{definition}[Specific Volume Quantity]
  \label{def:physics:phyx-mini-0463:phyxminiproblems-problemphyxmini0463-specificvolumequantity}
  \lean{PhyXMiniProblems.ProblemPhyXMini0463.SpecificVolumeQuantity}
  \uses{def:physics:phyx-mini-0463:phyxminiproblems-problemphyxmini0463-specificvolumedimension}
  A nonnegative, unit-independent thermodynamic specific volume
\end{definition}
\begin{proof}
  This is the declared model definition of Specific Volume Quantity.
\end{proof}
\begin{definition}[Specific Energy Quantity]
  \label{def:physics:phyx-mini-0463:phyxminiproblems-problemphyxmini0463-specificenergyquantity}
  \lean{PhyXMiniProblems.ProblemPhyXMini0463.SpecificEnergyQuantity}
  \uses{def:physics:phyx-mini-0463:phyxminiproblems-problemphyxmini0463-specificenergydimension}
  A signed, unit-independent thermodynamic specific internal energy
\end{definition}
\begin{proof}
  This is the declared model definition of Specific Energy Quantity.
\end{proof}
\begin{definition}[Acceleration Quantity]
  \label{def:physics:phyx-mini-0463:phyxminiproblems-problemphyxmini0463-accelerationquantity}
  \lean{PhyXMiniProblems.ProblemPhyXMini0463.AccelerationQuantity}
  \uses{def:physics:phyx-mini-0463:phyxminiproblems-problemphyxmini0463-accelerationdimension}
  A nonnegative acceleration magnitude
\end{definition}
\begin{proof}
  This is the declared model definition of Acceleration Quantity.
\end{proof}
\begin{definition}[Force Quantity]
  \label{def:physics:phyx-mini-0463:phyxminiproblems-problemphyxmini0463-forcequantity}
  \lean{PhyXMiniProblems.ProblemPhyXMini0463.ForceQuantity}
  \uses{def:physics:phyx-mini-0463:phyxminiproblems-problemphyxmini0463-forcedimension}
  A nonnegative force magnitude
\end{definition}
\begin{proof}
  This is the declared model definition of Force Quantity.
\end{proof}
\begin{definition}[mass In Kilograms]
  \label{def:physics:phyx-mini-0463:phyxminiproblems-problemphyxmini0463-massinkilograms}
  \lean{PhyXMiniProblems.ProblemPhyXMini0463.massInKilograms}
  \uses{def:physics:phyx-mini-0463:phyxminiproblems-problemphyxmini0463-massquantity}
  Read a physical mass in kilograms
\end{definition}
\begin{proof}
  This is the declared model definition of mass In Kilograms.
\end{proof}
\begin{definition}[volume In Cubic Meters]
  \label{def:physics:phyx-mini-0463:phyxminiproblems-problemphyxmini0463-volumeincubicmeters}
  \lean{PhyXMiniProblems.ProblemPhyXMini0463.volumeInCubicMeters}
  \uses{def:physics:phyx-mini-0463:phyxminiproblems-problemphyxmini0463-volumequantity}
  Read a physical volume in cubic metres
\end{definition}
\begin{proof}
  This is the declared model definition of volume In Cubic Meters.
\end{proof}
\begin{definition}[specific Volume In Cubic Meters Per Kilogram]
  \label{def:physics:phyx-mini-0463:phyxminiproblems-problemphyxmini0463-specificvolumeincubicmetersperkilogram}
  \lean{PhyXMiniProblems.ProblemPhyXMini0463.specificVolumeInCubicMetersPerKilogram}
  \uses{def:physics:phyx-mini-0463:phyxminiproblems-problemphyxmini0463-specificvolumequantity}
  Read a specific volume in cubic metres per kilogram
\end{definition}
\begin{proof}
  This is the declared model definition of specific Volume In Cubic Meters Per Kilogram.
\end{proof}
\begin{definition}[area In Square Meters]
  \label{def:physics:phyx-mini-0463:phyxminiproblems-problemphyxmini0463-areainsquaremeters}
  \lean{PhyXMiniProblems.ProblemPhyXMini0463.areaInSquareMeters}
  Read a piston area in square metres
\end{definition}
\begin{proof}
  This is the declared model definition of area In Square Meters.
\end{proof}
\begin{definition}[pressure In Pascals]
  \label{def:physics:phyx-mini-0463:phyxminiproblems-problemphyxmini0463-pressureinpascals}
  \lean{PhyXMiniProblems.ProblemPhyXMini0463.pressureInPascals}
  Read a pressure in pascals
\end{definition}
\begin{proof}
  This is the declared model definition of pressure In Pascals.
\end{proof}
\begin{definition}[pressure In Kilopascals]
  \label{def:physics:phyx-mini-0463:phyxminiproblems-problemphyxmini0463-pressureinkilopascals}
  \lean{PhyXMiniProblems.ProblemPhyXMini0463.pressureInKilopascals}
  \uses{def:physics:phyx-mini-0463:phyxminiproblems-problemphyxmini0463-pressureinpascals}
  Read a pressure in kilopascals
\end{definition}
\begin{proof}
  This is the declared model definition of pressure In Kilopascals.
\end{proof}
\begin{definition}[acceleration In Meters Per Second Squared]
  \label{def:physics:phyx-mini-0463:phyxminiproblems-problemphyxmini0463-accelerationinmeterspersecondsquared}
  \lean{PhyXMiniProblems.ProblemPhyXMini0463.accelerationInMetersPerSecondSquared}
  \uses{def:physics:phyx-mini-0463:phyxminiproblems-problemphyxmini0463-accelerationquantity}
  Read an acceleration magnitude in metres per second squared
\end{definition}
\begin{proof}
  This is the declared model definition of acceleration In Meters Per Second Squared.
\end{proof}
\begin{definition}[force In Newtons]
  \label{def:physics:phyx-mini-0463:phyxminiproblems-problemphyxmini0463-forceinnewtons}
  \lean{PhyXMiniProblems.ProblemPhyXMini0463.forceInNewtons}
  \uses{def:physics:phyx-mini-0463:phyxminiproblems-problemphyxmini0463-forcequantity}
  Read a force magnitude in newtons
\end{definition}
\begin{proof}
  This is the declared model definition of force In Newtons.
\end{proof}
\begin{definition}[energy In Joules]
  \label{def:physics:phyx-mini-0463:phyxminiproblems-problemphyxmini0463-energyinjoules}
  \lean{PhyXMiniProblems.ProblemPhyXMini0463.energyInJoules}
  Read heat or work in joules
\end{definition}
\begin{proof}
  This is the declared model definition of energy In Joules.
\end{proof}
\begin{definition}[energy In Kilojoules]
  \label{def:physics:phyx-mini-0463:phyxminiproblems-problemphyxmini0463-energyinkilojoules}
  \lean{PhyXMiniProblems.ProblemPhyXMini0463.energyInKilojoules}
  \uses{def:physics:phyx-mini-0463:phyxminiproblems-problemphyxmini0463-energyinjoules}
  Read heat or work in kilojoules
\end{definition}
\begin{proof}
  This is the declared model definition of energy In Kilojoules.
\end{proof}
\begin{definition}[specific Energy In Kilojoules Per Kilogram]
  \label{def:physics:phyx-mini-0463:phyxminiproblems-problemphyxmini0463-specificenergyinkilojoulesperkilogram}
  \lean{PhyXMiniProblems.ProblemPhyXMini0463.specificEnergyInKilojoulesPerKilogram}
  \uses{def:physics:phyx-mini-0463:phyxminiproblems-problemphyxmini0463-specificenergyquantity}
  Read specific internal energy in kilojoules per kilogram
\end{definition}
\begin{proof}
  This is the declared model definition of specific Energy In Kilojoules Per Kilogram.
\end{proof}
\begin{definition}[Process State]
  \label{def:physics:phyx-mini-0463:phyxminiproblems-problemphyxmini0463-processstate}
  \lean{PhyXMiniProblems.ProblemPhyXMini0463.ProcessState}
  Distinguished equilibrium states of the two-stage heating process
\end{definition}
\begin{proof}
  This is the declared model definition of Process State.
\end{proof}
\begin{definition}[Fluid Substance]
  \label{def:physics:phyx-mini-0463:phyxminiproblems-problemphyxmini0463-fluidsubstance}
  \lean{PhyXMiniProblems.ProblemPhyXMini0463.FluidSubstance}
  Substance confined below the piston
\end{definition}
\begin{proof}
  This is the declared model definition of Fluid Substance.
\end{proof}
\begin{definition}[Water Phase Region]
  \label{def:physics:phyx-mini-0463:phyxminiproblems-problemphyxmini0463-waterphaseregion}
  \lean{PhyXMiniProblems.ProblemPhyXMini0463.WaterPhaseRegion}
  Thermodynamic phase region occupied by a water state
\end{definition}
\begin{proof}
  This is the declared model definition of Water Phase Region.
\end{proof}
\begin{definition}[Piston Support Regime]
  \label{def:physics:phyx-mini-0463:phyxminiproblems-problemphyxmini0463-pistonsupportregime}
  \lean{PhyXMiniProblems.ProblemPhyXMini0463.PistonSupportRegime}
  Mechanical support regime of the piston at a distinguished state
\end{definition}
\begin{proof}
  This is the declared model definition of Piston Support Regime.
\end{proof}
\begin{definition}[Heating Process Regime]
  \label{def:physics:phyx-mini-0463:phyxminiproblems-problemphyxmini0463-heatingprocessregime}
  \lean{PhyXMiniProblems.ProblemPhyXMini0463.HeatingProcessRegime}
  Idealized heating path selected by the stops and piston load
\end{definition}
\begin{proof}
  This is the declared model definition of Heating Process Regime.
\end{proof}
\begin{definition}[Heat Transfer Sign Convention]
  \label{def:physics:phyx-mini-0463:phyxminiproblems-problemphyxmini0463-heattransfersignconvention}
  \lean{PhyXMiniProblems.ProblemPhyXMini0463.HeatTransferSignConvention}
  Sign convention for the requested heat transfer
\end{definition}
\begin{proof}
  This is the declared model definition of Heat Transfer Sign Convention.
\end{proof}
\begin{definition}[Vertical Direction]
  \label{def:physics:phyx-mini-0463:phyxminiproblems-problemphyxmini0463-verticaldirection}
  \lean{PhyXMiniProblems.ProblemPhyXMini0463.VerticalDirection}
  Vertical directions visible in the supplied figure
\end{definition}
\begin{proof}
  This is the declared model definition of Vertical Direction.
\end{proof}
\begin{definition}[Water State]
  \label{def:physics:phyx-mini-0463:phyxminiproblems-problemphyxmini0463-waterstate}
  \lean{PhyXMiniProblems.ProblemPhyXMini0463.WaterState}
  \uses{def:physics:phyx-mini-0463:phyxminiproblems-problemphyxmini0463-volumequantity, def:physics:phyx-mini-0463:phyxminiproblems-problemphyxmini0463-specificvolumequantity, def:physics:phyx-mini-0463:phyxminiproblems-problemphyxmini0463-specificenergyquantity, def:physics:phyx-mini-0463:phyxminiproblems-problemphyxmini0463-waterphaseregion}
  One equilibrium state of the closed water sample
\end{definition}
\begin{proof}
  This is the declared model definition of Water State.
\end{proof}
\begin{definition}[Water Table Interpolation Convention]
  \label{def:physics:phyx-mini-0463:phyxminiproblems-problemphyxmini0463-watertableinterpolationconvention}
  \lean{PhyXMiniProblems.ProblemPhyXMini0463.WaterTableInterpolationConvention}
  How an off-grid state is interpolated in the external steam table
\end{definition}
\begin{proof}
  This is the declared model definition of Water Table Interpolation Convention.
\end{proof}
\begin{definition}[Water Property Table]
  \label{def:physics:phyx-mini-0463:phyxminiproblems-problemphyxmini0463-waterpropertytable}
  \lean{PhyXMiniProblems.ProblemPhyXMini0463.WaterPropertyTable}
  \uses{def:physics:phyx-mini-0463:phyxminiproblems-problemphyxmini0463-specificvolumequantity, def:physics:phyx-mini-0463:phyxminiproblems-problemphyxmini0463-specificenergyquantity, def:physics:phyx-mini-0463:phyxminiproblems-problemphyxmini0463-watertableinterpolationconvention}
  An abstract equilibrium water-property table. Saturated-liquid and saturated-vapor entries are indexed by pressure; a single-phase state can be looked up by pressure and specific volume. These are constitutive properties, not definitions of the requested heat transfer. The source does not identify a steam-table edition. The identifier is therefore retained as a parameter rather than silently fixed here
\end{definition}
\begin{proof}
  This is the declared model definition of Water Property Table.
\end{proof}
\begin{definition}[Figure Object]
  \label{def:physics:phyx-mini-0463:phyxminiproblems-problemphyxmini0463-figureobject}
  \lean{PhyXMiniProblems.ProblemPhyXMini0463.FigureObject}
  Physical objects visible in the supplied raster image
\end{definition}
\begin{proof}
  This is the declared model definition of Figure Object.
\end{proof}
\begin{definition}[Figure Label]
  \label{def:physics:phyx-mini-0463:phyxminiproblems-problemphyxmini0463-figurelabel}
  \lean{PhyXMiniProblems.ProblemPhyXMini0463.FigureLabel}
  Literal labels visible in the supplied raster image
\end{definition}
\begin{proof}
  This is the declared model definition of Figure Label.
\end{proof}
\begin{definition}[Stopped Piston Cylinder Figure]
  \label{def:physics:phyx-mini-0463:phyxminiproblems-problemphyxmini0463-stoppedpistoncylinderfigure}
  \lean{PhyXMiniProblems.ProblemPhyXMini0463.StoppedPistonCylinderFigure}
  \uses{def:physics:phyx-mini-0463:phyxminiproblems-problemphyxmini0463-verticaldirection, def:physics:phyx-mini-0463:phyxminiproblems-problemphyxmini0463-figureobject, def:physics:phyx-mini-0463:phyxminiproblems-problemphyxmini0463-figurelabel}
  Qualitative evidence transcribed from the primary image. The image contains no numerical scale or thermodynamic-state readout
\end{definition}
\begin{proof}
  This is the declared model definition of Stopped Piston Cylinder Figure.
\end{proof}
\begin{definition}[Stopped Piston Water Setup]
  \label{def:physics:phyx-mini-0463:phyxminiproblems-problemphyxmini0463-stoppedpistonwatersetup}
  \lean{PhyXMiniProblems.ProblemPhyXMini0463.StoppedPistonWaterSetup}
  \uses{def:physics:phyx-mini-0463:phyxminiproblems-problemphyxmini0463-massquantity, def:physics:phyx-mini-0463:phyxminiproblems-problemphyxmini0463-accelerationquantity, def:physics:phyx-mini-0463:phyxminiproblems-problemphyxmini0463-forcequantity, def:physics:phyx-mini-0463:phyxminiproblems-problemphyxmini0463-processstate, def:physics:phyx-mini-0463:phyxminiproblems-problemphyxmini0463-fluidsubstance, def:physics:phyx-mini-0463:phyxminiproblems-problemphyxmini0463-pistonsupportregime, def:physics:phyx-mini-0463:phyxminiproblems-problemphyxmini0463-heatingprocessregime, def:physics:phyx-mini-0463:phyxminiproblems-problemphyxmini0463-heattransfersignconvention, def:physics:phyx-mini-0463:phyxminiproblems-problemphyxmini0463-waterstate, def:physics:phyx-mini-0463:phyxminiproblems-problemphyxmini0463-waterpropertytable, def:physics:phyx-mini-0463:phyxminiproblems-problemphyxmini0463-stoppedpistoncylinderfigure}
  The complete physical setup. Heat and boundary work are independent physical observables; neither is defined from an answer choice or from the desired numerical conclusion
\end{definition}
\begin{proof}
  This is the declared model definition of Stopped Piston Water Setup.
\end{proof}
\begin{definition}[Matches Stopped Piston Water Scenario]
  \label{def:physics:phyx-mini-0463:phyxminiproblems-problemphyxmini0463-matchesstoppedpistonwaterscenario}
  \lean{PhyXMiniProblems.ProblemPhyXMini0463.MatchesStoppedPistonWaterScenario}
  \uses{def:physics:phyx-mini-0463:phyxminiproblems-problemphyxmini0463-stoppedpistonwatersetup}
  Qualitative conditions stated or conventionally idealized in the problem
\end{definition}
\begin{proof}
  This is the declared model definition of Matches Stopped Piston Water Scenario.
\end{proof}
\begin{definition}[Matches Problem Readouts]
  \label{def:physics:phyx-mini-0463:phyxminiproblems-problemphyxmini0463-matchesproblemreadouts}
  \lean{PhyXMiniProblems.ProblemPhyXMini0463.MatchesProblemReadouts}
  \uses{def:physics:phyx-mini-0463:phyxminiproblems-problemphyxmini0463-massinkilograms, def:physics:phyx-mini-0463:phyxminiproblems-problemphyxmini0463-volumeincubicmeters, def:physics:phyx-mini-0463:phyxminiproblems-problemphyxmini0463-pressureinkilopascals, def:physics:phyx-mini-0463:phyxminiproblems-problemphyxmini0463-stoppedpistonwatersetup}
  Numerical information stated in the prose. The triple-volume condition is an endpoint datum; neither the requested heat nor an answer label occurs here
\end{definition}
\begin{proof}
  This is the declared model definition of Matches Problem Readouts.
\end{proof}
\begin{definition}[Matches Supplied Stopped Piston Figure]
  \label{def:physics:phyx-mini-0463:phyxminiproblems-problemphyxmini0463-matchessuppliedstoppedpistonfigure}
  \lean{PhyXMiniProblems.ProblemPhyXMini0463.MatchesSuppliedStoppedPistonFigure}
  \uses{def:physics:phyx-mini-0463:phyxminiproblems-problemphyxmini0463-stoppedpistonwatersetup}
  Exact objects, labels, and geometric relations visible in the image
\end{definition}
\begin{proof}
  This is the declared model definition of Matches Supplied Stopped Piston Figure.
\end{proof}
\begin{definition}[Has Physical Stopped Piston Parameters]
  \label{def:physics:phyx-mini-0463:phyxminiproblems-problemphyxmini0463-hasphysicalstoppedpistonparameters}
  \lean{PhyXMiniProblems.ProblemPhyXMini0463.HasPhysicalStoppedPistonParameters}
  \uses{def:physics:phyx-mini-0463:phyxminiproblems-problemphyxmini0463-massinkilograms, def:physics:phyx-mini-0463:phyxminiproblems-problemphyxmini0463-volumeincubicmeters, def:physics:phyx-mini-0463:phyxminiproblems-problemphyxmini0463-specificvolumeincubicmetersperkilogram, def:physics:phyx-mini-0463:phyxminiproblems-problemphyxmini0463-areainsquaremeters, def:physics:phyx-mini-0463:phyxminiproblems-problemphyxmini0463-pressureinpascals, def:physics:phyx-mini-0463:phyxminiproblems-problemphyxmini0463-accelerationinmeterspersecondsquared, def:physics:phyx-mini-0463:phyxminiproblems-problemphyxmini0463-stoppedpistonwatersetup}
  Positivity and quality bounds selecting the physical branch of the model
\end{definition}
\begin{proof}
  This is the declared model definition of Has Physical Stopped Piston Parameters.
\end{proof}
\begin{definition}[Satisfies Water Property Relations]
  \label{def:physics:phyx-mini-0463:phyxminiproblems-problemphyxmini0463-satisfieswaterpropertyrelations}
  \lean{PhyXMiniProblems.ProblemPhyXMini0463.SatisfiesWaterPropertyRelations}
  \uses{def:physics:phyx-mini-0463:phyxminiproblems-problemphyxmini0463-specificvolumeincubicmetersperkilogram, def:physics:phyx-mini-0463:phyxminiproblems-problemphyxmini0463-specificenergyinkilojoulesperkilogram, def:physics:phyx-mini-0463:phyxminiproblems-problemphyxmini0463-stoppedpistonwatersetup}
  Constitutive state relations: saturated-mixture interpolation at the initial state and a pressure/specific-volume lookup at the final state. The equations contain no heat-transfer value
\end{definition}
\begin{proof}
  This is the declared model definition of Satisfies Water Property Relations.
\end{proof}
\begin{definition}[Matches Exercise Water Property Data]
  \label{def:physics:phyx-mini-0463:phyxminiproblems-problemphyxmini0463-matchesexercisewaterpropertydata}
  \lean{PhyXMiniProblems.ProblemPhyXMini0463.MatchesExerciseWaterPropertyData}
  \uses{def:physics:phyx-mini-0463:phyxminiproblems-problemphyxmini0463-specificvolumeincubicmetersperkilogram, def:physics:phyx-mini-0463:phyxminiproblems-problemphyxmini0463-specificenergyinkilojoulesperkilogram, def:physics:phyx-mini-0463:phyxminiproblems-problemphyxmini0463-stoppedpistonwatersetup}
  Calibrated equilibrium-water data used by the numerical exercise. The final entry is a linearly interpolated superheated-water readout, recorded to the nearest 0.1 kJ/kg; hence its exact value lies within 0.05 kJ/kg of 3718.8 kJ/kg. The table edition remains unconstrained because it is absent from the source. These constitutive readouts mention neither total heat nor any displayed answer choice
\end{definition}
\begin{proof}
  This is the declared model definition of Matches Exercise Water Property Data.
\end{proof}
\begin{definition}[Satisfies Closed Water Mass Volume Law]
  \label{def:physics:phyx-mini-0463:phyxminiproblems-problemphyxmini0463-satisfiesclosedwatermassvolumelaw}
  \lean{PhyXMiniProblems.ProblemPhyXMini0463.SatisfiesClosedWaterMassVolumeLaw}
  \uses{def:physics:phyx-mini-0463:phyxminiproblems-problemphyxmini0463-massinkilograms, def:physics:phyx-mini-0463:phyxminiproblems-problemphyxmini0463-volumeincubicmeters, def:physics:phyx-mini-0463:phyxminiproblems-problemphyxmini0463-specificvolumeincubicmetersperkilogram, def:physics:phyx-mini-0463:phyxminiproblems-problemphyxmini0463-stoppedpistonwatersetup}
  Conservation of water mass and V = m v at each modeled equilibrium state
\end{definition}
\begin{proof}
  This is the declared model definition of Satisfies Closed Water Mass Volume Law.
\end{proof}
\begin{definition}[Satisfies Piston Load Balance]
  \label{def:physics:phyx-mini-0463:phyxminiproblems-problemphyxmini0463-satisfiespistonloadbalance}
  \lean{PhyXMiniProblems.ProblemPhyXMini0463.SatisfiesPistonLoadBalance}
  \uses{def:physics:phyx-mini-0463:phyxminiproblems-problemphyxmini0463-massinkilograms, def:physics:phyx-mini-0463:phyxminiproblems-problemphyxmini0463-areainsquaremeters, def:physics:phyx-mini-0463:phyxminiproblems-problemphyxmini0463-pressureinpascals, def:physics:phyx-mini-0463:phyxminiproblems-problemphyxmini0463-accelerationinmeterspersecondsquared, def:physics:phyx-mini-0463:phyxminiproblems-problemphyxmini0463-forceinnewtons, def:physics:phyx-mini-0463:phyxminiproblems-problemphyxmini0463-stoppedpistonwatersetup}
  Mechanical equilibrium of the piston. At lift-off the lower-stop reaction is zero, so the water-pressure force balances the outside-pressure force and piston weight. The initial equation records the nonzero support reaction
\end{definition}
\begin{proof}
  This is the declared model definition of Satisfies Piston Load Balance.
\end{proof}
\begin{definition}[Satisfies Stopped Then Constant Pressure Path]
  \label{def:physics:phyx-mini-0463:phyxminiproblems-problemphyxmini0463-satisfiesstoppedthenconstantpressurepath}
  \lean{PhyXMiniProblems.ProblemPhyXMini0463.SatisfiesStoppedThenConstantPressurePath}
  \uses{def:physics:phyx-mini-0463:phyxminiproblems-problemphyxmini0463-stoppedpistonwatersetup}
  The stops enforce constant volume before lift-off. Once the piston floats, the fixed outside pressure and piston weight impose the constant float pressure through the final equilibrium state
\end{definition}
\begin{proof}
  This is the declared model definition of Satisfies Stopped Then Constant Pressure Path.
\end{proof}
\begin{definition}[Satisfies Two Stage Boundary Work Law]
  \label{def:physics:phyx-mini-0463:phyxminiproblems-problemphyxmini0463-satisfiestwostageboundaryworklaw}
  \lean{PhyXMiniProblems.ProblemPhyXMini0463.SatisfiesTwoStageBoundaryWorkLaw}
  \uses{def:physics:phyx-mini-0463:phyxminiproblems-problemphyxmini0463-volumeincubicmeters, def:physics:phyx-mini-0463:phyxminiproblems-problemphyxmini0463-pressureinkilopascals, def:physics:phyx-mini-0463:phyxminiproblems-problemphyxmini0463-energyinkilojoules, def:physics:phyx-mini-0463:phyxminiproblems-problemphyxmini0463-stoppedpistonwatersetup}
  Boundary work is zero in the stop-constrained stage and equals p ΔV in the subsequent quasistatic constant-pressure stage. Since kPa · m³ = kJ, the law is written directly in those calibrated readouts
\end{definition}
\begin{proof}
  This is the declared model definition of Satisfies Two Stage Boundary Work Law.
\end{proof}
\begin{definition}[Satisfies Closed System First Law]
  \label{def:physics:phyx-mini-0463:phyxminiproblems-problemphyxmini0463-satisfiesclosedsystemfirstlaw}
  \lean{PhyXMiniProblems.ProblemPhyXMini0463.SatisfiesClosedSystemFirstLaw}
  \uses{def:physics:phyx-mini-0463:phyxminiproblems-problemphyxmini0463-massinkilograms, def:physics:phyx-mini-0463:phyxminiproblems-problemphyxmini0463-energyinkilojoules, def:physics:phyx-mini-0463:phyxminiproblems-problemphyxmini0463-specificenergyinkilojoulesperkilogram, def:physics:phyx-mini-0463:phyxminiproblems-problemphyxmini0463-stoppedpistonwatersetup}
  The closed-system first law with the sign conventions Q > 0 into the water and W > 0 done by the water. Changes in kinetic and potential energy are neglected as stated in the scenario assumptions
\end{definition}
\begin{proof}
  This is the declared model definition of Satisfies Closed System First Law.
\end{proof}
\begin{definition}[Answer Choice]
  \label{def:physics:phyx-mini-0463:phyxminiproblems-problemphyxmini0463-answerchoice}
  \lean{PhyXMiniProblems.ProblemPhyXMini0463.AnswerChoice}
  Labels printed beside the four candidate heat-transfer values
\end{definition}
\begin{proof}
  This is the declared model definition of Answer Choice.
\end{proof}
\begin{definition}[displayed Heat Transfer In Kilojoules]
  \label{def:physics:phyx-mini-0463:phyxminiproblems-problemphyxmini0463-displayedheattransferinkilojoules}
  \lean{PhyXMiniProblems.ProblemPhyXMini0463.displayedHeatTransferInKilojoules}
  \uses{def:physics:phyx-mini-0463:phyxminiproblems-problemphyxmini0463-answerchoice}
  Heat transfer in kilojoules printed beside each answer label
\end{definition}
\begin{proof}
  This is the declared model definition of displayed Heat Transfer In Kilojoules.
\end{proof}
\begin{definition}[recorded Dataset Answer]
  \label{def:physics:phyx-mini-0463:phyxminiproblems-problemphyxmini0463-recordeddatasetanswer}
  \lean{PhyXMiniProblems.ProblemPhyXMini0463.recordedDatasetAnswer}
  \uses{def:physics:phyx-mini-0463:phyxminiproblems-problemphyxmini0463-answerchoice}
  Dataset answer metadata, deliberately not used as a theorem premise
\end{definition}
\begin{proof}
  This is the declared model definition of recorded Dataset Answer.
\end{proof}
\begin{definition}[heat Transfer Choice Error]
  \label{def:physics:phyx-mini-0463:phyxminiproblems-problemphyxmini0463-heattransferchoiceerror}
  \lean{PhyXMiniProblems.ProblemPhyXMini0463.heatTransferChoiceError}
  \uses{def:physics:phyx-mini-0463:phyxminiproblems-problemphyxmini0463-energyinkilojoules, def:physics:phyx-mini-0463:phyxminiproblems-problemphyxmini0463-stoppedpistonwatersetup, def:physics:phyx-mini-0463:phyxminiproblems-problemphyxmini0463-answerchoice, def:physics:phyx-mini-0463:phyxminiproblems-problemphyxmini0463-displayedheattransferinkilojoules}
  Absolute reporting error, in kilojoules, for one displayed choice
\end{definition}
\begin{proof}
  This is the declared model definition of heat Transfer Choice Error.
\end{proof}
\begin{definition}[Is Closest Reported Heat Transfer Choice]
  \label{def:physics:phyx-mini-0463:phyxminiproblems-problemphyxmini0463-isclosestreportedheattransferchoice}
  \lean{PhyXMiniProblems.ProblemPhyXMini0463.IsClosestReportedHeatTransferChoice}
  \uses{def:physics:phyx-mini-0463:phyxminiproblems-problemphyxmini0463-stoppedpistonwatersetup, def:physics:phyx-mini-0463:phyxminiproblems-problemphyxmini0463-answerchoice, def:physics:phyx-mini-0463:phyxminiproblems-problemphyxmini0463-heattransferchoiceerror}
  A displayed choice is at least as close as every other displayed value
\end{definition}
\begin{proof}
  This is the declared model definition of Is Closest Reported Heat Transfer Choice.
\end{proof}
\begin{definition}[Is Unique Closest Reported Heat Transfer Choice]
  \label{def:physics:phyx-mini-0463:phyxminiproblems-problemphyxmini0463-isuniqueclosestreportedheattransferchoice}
  \lean{PhyXMiniProblems.ProblemPhyXMini0463.IsUniqueClosestReportedHeatTransferChoice}
  \uses{def:physics:phyx-mini-0463:phyxminiproblems-problemphyxmini0463-stoppedpistonwatersetup, def:physics:phyx-mini-0463:phyxminiproblems-problemphyxmini0463-answerchoice, def:physics:phyx-mini-0463:phyxminiproblems-problemphyxmini0463-isclosestreportedheattransferchoice}
  A choice is the unique closest displayed value to the derived heat
\end{definition}
\begin{proof}
  This is the declared model definition of Is Unique Closest Reported Heat Transfer Choice.
\end{proof}
\begin{lemma}[initial Mixture Specific Properties]
  \label{lem:physics:phyx-mini-0463:phyxminiproblems-problemphyxmini0463-initialmixturespecificproperties}
  \lean{PhyXMiniProblems.ProblemPhyXMini0463.initialMixtureSpecificProperties}
  \uses{def:physics:phyx-mini-0463:phyxminiproblems-problemphyxmini0463-specificvolumeincubicmetersperkilogram, def:physics:phyx-mini-0463:phyxminiproblems-problemphyxmini0463-specificenergyinkilojoulesperkilogram, def:physics:phyx-mini-0463:phyxminiproblems-problemphyxmini0463-stoppedpistonwatersetup, def:physics:phyx-mini-0463:phyxminiproblems-problemphyxmini0463-matchesproblemreadouts, def:physics:phyx-mini-0463:phyxminiproblems-problemphyxmini0463-satisfieswaterpropertyrelations, def:physics:phyx-mini-0463:phyxminiproblems-problemphyxmini0463-matchesexercisewaterpropertydata}
  The saturated-mixture interpolation gives the initial specific volume 0.8475215 m³/kg and specific internal energy 1461.73 kJ/kg under the exercise calibration
\end{lemma}
\begin{proof}
  The conclusion follows from the governing relations and figure readouts named in the statement of initial Mixture Specific Properties.
\end{proof}
\begin{lemma}[distinguished Volumes In Cubic Meters]
  \label{lem:physics:phyx-mini-0463:phyxminiproblems-problemphyxmini0463-distinguishedvolumesincubicmeters}
  \lean{PhyXMiniProblems.ProblemPhyXMini0463.distinguishedVolumesInCubicMeters}
  \uses{def:physics:phyx-mini-0463:phyxminiproblems-problemphyxmini0463-volumeincubicmeters, def:physics:phyx-mini-0463:phyxminiproblems-problemphyxmini0463-stoppedpistonwatersetup, def:physics:phyx-mini-0463:phyxminiproblems-problemphyxmini0463-matchesproblemreadouts, def:physics:phyx-mini-0463:phyxminiproblems-problemphyxmini0463-satisfieswaterpropertyrelations, def:physics:phyx-mini-0463:phyxminiproblems-problemphyxmini0463-matchesexercisewaterpropertydata, def:physics:phyx-mini-0463:phyxminiproblems-problemphyxmini0463-satisfiesclosedwatermassvolumelaw, def:physics:phyx-mini-0463:phyxminiproblems-problemphyxmini0463-satisfiesstoppedthenconstantpressurepath}
  Mass conservation and the triple-volume endpoint give V₁ = V₂ = 8.475215 m³ and V₃ = 25.425645 m³
\end{lemma}
\begin{proof}
  The conclusion follows from the governing relations and figure readouts named in the statement of distinguished Volumes In Cubic Meters.
\end{proof}
\begin{lemma}[boundary Work In Kilojoules]
  \label{lem:physics:phyx-mini-0463:phyxminiproblems-problemphyxmini0463-boundaryworkinkilojoules}
  \lean{PhyXMiniProblems.ProblemPhyXMini0463.boundaryWorkInKilojoules}
  \uses{def:physics:phyx-mini-0463:phyxminiproblems-problemphyxmini0463-energyinkilojoules, def:physics:phyx-mini-0463:phyxminiproblems-problemphyxmini0463-stoppedpistonwatersetup, def:physics:phyx-mini-0463:phyxminiproblems-problemphyxmini0463-matchesproblemreadouts, def:physics:phyx-mini-0463:phyxminiproblems-problemphyxmini0463-satisfieswaterpropertyrelations, def:physics:phyx-mini-0463:phyxminiproblems-problemphyxmini0463-matchesexercisewaterpropertydata, def:physics:phyx-mini-0463:phyxminiproblems-problemphyxmini0463-satisfiesclosedwatermassvolumelaw, def:physics:phyx-mini-0463:phyxminiproblems-problemphyxmini0463-satisfiesstoppedthenconstantpressurepath, def:physics:phyx-mini-0463:phyxminiproblems-problemphyxmini0463-satisfiestwostageboundaryworklaw}
  Only the moving-piston stage performs boundary work; its constant-pressure value is 3390.086 kJ. This is answer choice C's number, but it is work, not the requested heat transfer
\end{lemma}
\begin{proof}
  The conclusion follows from the governing relations and figure readouts named in the statement of boundary Work In Kilojoules.
\end{proof}
\begin{lemma}[heat Transfer Estimate In Kilojoules]
  \label{lem:physics:phyx-mini-0463:phyxminiproblems-problemphyxmini0463-heattransferestimateinkilojoules}
  \lean{PhyXMiniProblems.ProblemPhyXMini0463.heatTransferEstimateInKilojoules}
  \uses{def:physics:phyx-mini-0463:phyxminiproblems-problemphyxmini0463-energyinkilojoules, def:physics:phyx-mini-0463:phyxminiproblems-problemphyxmini0463-stoppedpistonwatersetup, def:physics:phyx-mini-0463:phyxminiproblems-problemphyxmini0463-matchesproblemreadouts, def:physics:phyx-mini-0463:phyxminiproblems-problemphyxmini0463-satisfieswaterpropertyrelations, def:physics:phyx-mini-0463:phyxminiproblems-problemphyxmini0463-matchesexercisewaterpropertydata, def:physics:phyx-mini-0463:phyxminiproblems-problemphyxmini0463-satisfiesclosedwatermassvolumelaw, def:physics:phyx-mini-0463:phyxminiproblems-problemphyxmini0463-satisfiesstoppedthenconstantpressurepath, def:physics:phyx-mini-0463:phyxminiproblems-problemphyxmini0463-satisfiestwostageboundaryworklaw, def:physics:phyx-mini-0463:phyxminiproblems-problemphyxmini0463-satisfiesclosedsystemfirstlaw}
  Combining the independently calibrated endpoint water properties, boundary work, and the first law places the heat within 0.5 kJ of 25960.786 kJ. This retains the precision of the interpolated property readout instead of reverse-calibrating it to an exact answer choice
\end{lemma}
\begin{proof}
  The conclusion follows from the governing relations and figure readouts named in the statement of heat Transfer Estimate In Kilojoules.
\end{proof}
% --- Archon declaration topology end ---
% --- Archon physics formalization source end ---
